\documentclass[a4paper,11pt]{article}

\usepackage{fullpage}
\usepackage[utf8]{inputenc}
\usepackage[T1]{fontenc}
\usepackage{amsmath,amsthm}
\usepackage{amsfonts}
\usepackage{graphicx}
\usepackage{latexsym}
\usepackage{float}
\usepackage{comment}
\usepackage{array}
\usepackage{booktabs}

% --- Packages ajoutés pour l'enrichissement pédagogique ---
\usepackage{xcolor}
\usepackage{listings}
\usepackage[most]{tcolorbox}
\usepackage{hyperref}

%%%%%%%%%%%%%%% CONFIGURATION LISTINGS %%%%%%%%%%%%%%%%%%%

\definecolor{codegreen}{rgb}{0.25,0.49,0.25}
\definecolor{codegray}{rgb}{0.5,0.5,0.5}
\definecolor{codepurple}{rgb}{0.58,0,0.82}
\definecolor{codered}{rgb}{0.7,0.1,0.1}
\definecolor{codeblue}{rgb}{0.0,0.0,0.7}
\definecolor{backcolour}{rgb}{0.97,0.97,0.97}

\lstset{
    language=C,
    backgroundcolor=\color{backcolour},
    commentstyle=\color{codegreen}\itshape,
    keywordstyle=\color{codeblue}\bfseries,
    stringstyle=\color{codered},
    basicstyle=\ttfamily\small,
    breakatwhitespace=false,
    breaklines=true,
    captionpos=b,
    keepspaces=true,
    numbers=left,
    numbersep=8pt,
    numberstyle=\tiny\color{codegray},
    showspaces=false,
    showstringspaces=false,
    showtabs=false,
    tabsize=4,
    frame=single,
    rulecolor=\color{codegray},
    xleftmargin=1.5em,
    framexleftmargin=1.5em,
    literate={é}{{\'e}}1 {è}{{\`e}}1 {ê}{{\^e}}1 {à}{{\`a}}1 {ù}{{\`u}}1 {î}{{\^i}}1 {ô}{{\^o}}1
}

%%%%%%%%%%%%%%% ENVIRONNEMENTS TCOLORBOX %%%%%%%%%%%%%%%%%%

% Rappel de cours (fond bleu clair)
\newtcolorbox{rappel}[1][]{
    colback=blue!5!white,
    colframe=blue!50!black,
    fonttitle=\bfseries,
    title={\large Rappel de cours},
    #1
}

% Point clé (fond vert clair)
\newtcolorbox{pointcle}[1][]{
    colback=green!5!white,
    colframe=green!40!black,
    fonttitle=\bfseries,
    title={Point clé à retenir},
    #1
}

% Attention (fond orange clair)
\newtcolorbox{attention}[1][]{
    colback=orange!8!white,
    colframe=orange!60!black,
    fonttitle=\bfseries,
    title={Attention},
    #1
}

% Avertissement éthique (fond rouge)
\newtcolorbox{ethique}[1][]{
    colback=red!5!white,
    colframe=red!60!black,
    fonttitle=\bfseries\large,
    title={Avertissement éthique et légal},
    #1
}

%%%%%%%%%%%%%%% COMMANDES %%%%%%%%%%%%%%%%%%%%%%%%%%%

\newcommand{\itemb}{\item[$\bullet$]}

\newcommand{\Fd}{\mathbb{F}}
\newcommand{\Znat}{\mathbb{Z}}
\newcommand{\Nnat}{\mathbb{N}}

%%%%%%%%%%%%%%% ENVIRONEMENTS %%%%%%%%%%%%%%%%%%%%%%%

\theoremstyle{plain}
\newtheorem{lemme}{Lemme}
\newtheorem{algorithme}{Algorithme}
\newtheorem{corolaire}{Corollaire}
\newtheorem{propriete}{Propri\'et\'e}
\newtheorem{proposition}{Proposition}
\newtheorem{theoreme}{Th\'eor\`eme}
\newtheorem{definition}{D\'efinition}

\theoremstyle{definition}
\newtheorem{exercice}{Exercice}
\newtheorem{remarque}{Remarque}
\newtheorem{exemple}{Exemple}

\hypersetup{
    colorlinks=true,
    linkcolor=blue!60!black,
    urlcolor=blue!60!black
}

%%%%%%%%%%%%%% DOCUMENTS %%%%%%%%%%%%%%%%%%%%%%%%%%%%


\begin{document}

\hbox{\raisebox{0.4em}{\vrule depth 0pt height 0.5pt width \textwidth}}
\noindent \textbf{Université de Perpignan} \hfill  L3 Informatique \\
 \hfill C. Legrand \hfill  Année \the\year \\
\smallskip
\begin{center}
{
  {\scshape \large TD5 Sécurité}  \\
 Techniques d'intrusion -- Débordement de tampon
}
\end{center}
\hbox{\raisebox{0.4em}{\vrule depth 0pt height 0.9pt width \textwidth}}

\bigskip

% --- Avertissement éthique ---
\begin{ethique}
Les techniques présentées dans ce TD sont étudiées dans un \textbf{cadre strictement académique} afin de comprendre les mécanismes utilisés par les attaquants et ainsi pouvoir mieux s'en protéger.

\textbf{Toute utilisation de ces techniques en dehors du cadre pédagogique est illégale} et passible de sanctions pénales (chapitre III du Code pénal, articles 323-1 à 323-8 -- \textit{Des atteintes aux systèmes de traitement automatisé de données}) :
\begin{itemize}
    \item Accès ou maintien frauduleux dans un STAD (art.~323-1) : \textbf{3 ans} d'emprisonnement et \textbf{100\,000\,\texteuro{}} d'amende
    \item Entrave au fonctionnement d'un STAD (art.~323-2) : \textbf{5 ans} d'emprisonnement et \textbf{150\,000\,\texteuro{}} d'amende
    \item Introduction frauduleuse de données dans un STAD (art.~323-3) : \textbf{5 ans} d'emprisonnement et \textbf{150\,000\,\texteuro{}} d'amende
\end{itemize}
Peines portées à \textbf{7 ans} et \textbf{300\,000\,\texteuro{}} lorsque le système visé traite des données à caractère personnel mis en \oe uvre par l'État, et jusqu'à \textbf{10 ans} et \textbf{300\,000\,\texteuro{}} en bande organisée (art.~323-4-1). \textit{Peines mises à jour par la loi LOPMI n\textsuperscript{o}2023-22 du 24 janvier 2023.}
\end{ethique}

\bigskip

\begin{tcolorbox}[colback=gray!5!white, colframe=gray!60!black, title={Options de compilation}]
Les options suivantes sont à utiliser lorsque l'on veut déborder sur la pile (\texttt{-fno-stack-protector}) et connaître les adresses du code (\texttt{-no-pie}) qui ne sont plus randomisées :
\begin{lstlisting}[language=bash,numbers=none]
gcc -fno-stack-protector -no-pie -g code.c
\end{lstlisting}
Une dernière option que l'on n'utilisera pas : si l'on souhaite exécuter du code injecté dans la pile (\texttt{-z execstack}).
\end{tcolorbox}

\medskip

\begin{rappel}
\textbf{Organisation de la pile (stack) :} lors de l'appel d'une fonction, le processeur empile :
\begin{enumerate}
    \item Les paramètres de la fonction (en partie via registres en x86-64)
    \item L'adresse de retour \texttt{savedRIP} (empilée par \texttt{call})
    \item L'ancien \texttt{RBP} (\texttt{savedRBP}) sauvegardé par \texttt{push rbp}
    \item Les variables locales de la fonction
\end{enumerate}
Le registre \texttt{RSP} pointe sur le sommet de la pile ; \texttt{RBP} pointe sur la base du cadre de pile courant.
\end{rappel}

\bigskip


%%%%%%%%%%%%%%%%%%%%%%%%%%%%%%%%%%%%%%%%%%%%%%%%%%%%%%%%%%%%%%%%%%
%% EXERCICE 1 : Débordement de tampon simple
%%%%%%%%%%%%%%%%%%%%%%%%%%%%%%%%%%%%%%%%%%%%%%%%%%%%%%%%%%%%%%%%%%

\begin{exercice}[Débordement de tampon]
\begin{enumerate}
\item
On considère le programme suivant :
\begin{lstlisting}[language=C]
#include <stdio.h>
#include <string.h>

int main(int argc, char *argv[]){
  int valeur=5;
  char tampon_un[8], tampon_deux[8];

  strcpy(tampon_un,"un");
  strcpy(tampon_deux,"deux");

  printf("[AVANT] tampon_deux est a %p et contient '%s'\n",
         tampon_deux, tampon_deux);
  printf("[AVANT] tampon_un est a %p et contient '%s'\n",
         tampon_un, tampon_un);
  printf("[AVANT] valeur est a %p et vaut %d (0x%08x)\n",
         &valeur, valeur, valeur);

  printf("\n [STRCPY] copie de %d octets dans tampon_deux\n\n",
         (int)strlen(argv[1]));
  strcpy(tampon_deux, argv[1]);

  printf("[APRES] tampon_deux est a %p et contient '%s'\n",
         tampon_deux, tampon_deux);
  printf("[APRES] tampon_un est a %p et contient '%s'\n",
         tampon_un, tampon_un);
  printf("[APRES] valeur est a %p et vaut %d (0x%08x)\n",
         &valeur, valeur, valeur);

  return(0);
}
\end{lstlisting}

En faisant varier l'argument, provoquer un débordement de tampon, analysez-le avec \texttt{gdb}.

\begin{pointcle}
La fonction \texttt{strcpy} ne vérifie \textbf{pas} la taille du tampon destination. Si la source est plus longue que le tampon, les octets en surplus écrasent les données adjacentes dans la pile.
\end{pointcle}

\item  On considère maintenant le bout de programme ci-dessous qui restreint les droits d'accès :

\begin{lstlisting}[language=C]
#include <stdio.h>
#include <stdlib.h>
#include <string.h>

int verif_authentication(char *password){
  int auth_flag = 0;
  char password_tampon[16];

  strcpy(password_tampon, password);

  if( strcmp(password_tampon, "brillig") == 0)
    auth_flag = 1;
  if(strcmp(password_tampon, "outgrabe") == 0)
    auth_flag = 1;

  return auth_flag;
}

int main( int argc, char *argv[]){
  if(argc < 2){
    printf("Usage: %s <password>\n", argv[0]);
    exit(0);
  }

  if(verif_authentication(argv[1])){
    printf("\n-=-=-=-=-=-=-=-=-=-=--=-=-\n");
    printf("      Acces accorde.\n");
    printf("-=-=-=-=-=-=-=-=-=-=--=-=-\n");
  }
  else
    printf("     Acces refuse\n");
}
\end{lstlisting}
Est-il possible de faire un débordement de tampon pour forcer l'accès (\texttt{auth\_flag} $\neq 0$) ? Que se passe-t-il si les variables sont déclarées dans l'ordre suivant :
\begin{lstlisting}[language=C,numbers=none]
  char password_buffer[16];
  int auth_flag = 0;
\end{lstlisting}
Peut-on contourner ce problème pour faire afficher quand même \texttt{Acces accorde} par débordement de tampon ? Expliquez en détail pourquoi ça fonctionne.
\end{enumerate}
\end{exercice}


%%%%%%%%%%%%%%%%%%%%%%%%%%%%%%%%%%%%%%%%%%%%%%%%%%%%%%%%%%%%%%%%%%
%% EXERCICE 2 : Modification d'adresse de retour (incrément)
%%%%%%%%%%%%%%%%%%%%%%%%%%%%%%%%%%%%%%%%%%%%%%%%%%%%%%%%%%%%%%%%%%

\begin{exercice}[Modification d'adresse de retour dans la pile]
Reprenez le code vu en cours (donné ci-dessous) et ajustez les valeurs de $a$ et de $b$ afin d'exécuter en retour de fonction :
\begin{itemize}
\item Directement la fin de la fonction \texttt{main} (donc pas de \texttt{printf("coucou1")} et \texttt{printf("coucou2")}).
  \item Que se passera-t-il si on revient à l'instruction \texttt{b=15} ?
\end{itemize}

\begin{lstlisting}[language=C]
#include <stdio.h>

void function(long a, long b)
{
  long buffer[2]={0,1};
  buffer[a]+=b;
}
void main()
{
  long a,b;
  a=1;
  b=15;
  function(a,b);
  printf("coucou1\n");
  printf("coucou2\n");
}
\end{lstlisting}

\begin{rappel}
Pour trouver les bonnes valeurs de \texttt{a} et \texttt{b}, il faut :
\begin{enumerate}
    \item Analyser le code assembleur avec \texttt{gdb} pour connaître la taille du cadre de pile de \texttt{function}
    \item Calculer le décalage entre \texttt{buffer[0]} et \texttt{savedRIP}
    \item Déterminer l'incrément nécessaire pour sauter par-dessus les instructions souhaitées
\end{enumerate}
\end{rappel}
\end{exercice}


%%%%%%%%%%%%%%%%%%%%%%%%%%%%%%%%%%%%%%%%%%%%%%%%%%%%%%%%%%%%%%%%%%
%% EXERCICE 3 : Écrasement d'adresse de retour (débordement)
%%%%%%%%%%%%%%%%%%%%%%%%%%%%%%%%%%%%%%%%%%%%%%%%%%%%%%%%%%%%%%%%%%

\begin{exercice}[Écrasement d'adresse de retour par débordement]
On considère le code suivant :
\begin{lstlisting}[language=C]
#include <stdio.h>
#include <stdlib.h>
#include <string.h>

void f(char *s)
{
  char nom[10];
  strcpy(nom,s);
  printf("\n");
  printf("%lx\n",*((unsigned long *)(nom+18)));
  printf("Bonjour %s, on est dans f\n",nom);
}

void secret()
{
  printf("on est dans secret\n");
}

int main(int argc, char **argv)
{
  printf("adresse de f=%p, et de secret=%p\n",f,secret);
  f(argv[1]);
}
\end{lstlisting}
En écrasant l'adresse de retour de l'appel à la fonction \texttt{f}, faire en sorte d'exécuter la fonction \texttt{secret}.

\begin{attention}
Pour passer en paramètre une chaîne de caractères spécifiés par leur code ASCII, on peut utiliser \texttt{python3} comme suit (pour les deux codes ASCII \texttt{0x11} et \texttt{0xff}) :
\begin{lstlisting}[language=bash,numbers=none]
$ python3 -c "print('\x11\xff')"
\end{lstlisting}
\end{attention}
\end{exercice}


%%%%%%%%%%%%%%%%%%%%%%%%%%%%%%%%%%%%%%%%%%%%%%%%%%%%%%%%%%%%%%%%%%
%% EXERCICE 4 : Chaîne de format printf
%%%%%%%%%%%%%%%%%%%%%%%%%%%%%%%%%%%%%%%%%%%%%%%%%%%%%%%%%%%%%%%%%%

\begin{exercice}[Exploitation de la chaîne de format \texttt{printf}]
On considère le code suivant compilé avec les options \texttt{-fno-stack-protector -no-pie} :
\begin{lstlisting}[language=C]
#include <stdio.h>
#include <stdlib.h>
#include <string.h>

int main(int argc, char *argv[]) {
  unsigned long d;
  char texte[512];

  static int test_var = 100*256*256+100*256+100;
  if(argc < 3) {
    printf("Usage: %s <text to print> <entier>\n", argv[0]);
    exit(0);
  }
  d=atoi(argv[2]);

  strcpy(texte, argv[1]);
  printf("La bonne facon d'afficher des entrees utilisateurs:\n");
  printf("%s", texte);
  printf("\nLa mauvaise facon d'afficher des entrees utilisateurs:\n");
  printf(texte);
  printf("\n");
  // Debug output
  printf("[*] test_var @ %p = %lu, %04x\n", &test_var, &test_var, test_var);
  printf("[*] adresse de d=%p, valeur de d=0x%016lx\n", &d, d);
  printf("[*] adresse de texte=%p\n", texte);
  exit(0);
}
\end{lstlisting}

\begin{enumerate}
\item Comprendre et tester le programme, puis appliquer ce que l'on a vu en cours pour afficher le contenu de la pile de \texttt{main} en choisissant bien l'argument du programme dans la ligne de commande.
\item Si l'argument du programme contient le format \texttt{"\%s"}, que va-t-il se passer ? Peut-on choisir l'adresse qui concerne \texttt{"\%s"} et l'appliquer pour lire la valeur de \texttt{test\_var} ?
\end{enumerate}

\begin{pointcle}
L'appel \texttt{printf(texte)} est dangereux car l'attaquant contrôle la chaîne de format. Il peut utiliser \texttt{\%x}, \texttt{\%lx}, \texttt{\%s} pour lire la mémoire, et même \texttt{\%n} pour écrire en mémoire !
\end{pointcle}
\end{exercice}


\end{document}
